\clearpage
\section{Casos de uso}

Los casos de uso describen las situaciones en las que un usuario interactúa
con \textit{Manualito} y la respuesta que debe ofrecer el sistema. Cada uno
concreta el objetivo de la interacción, sus condiciones y el resultado esperado,
y muestra la relación con los requisitos que desarrolla.

\subsection{Actores y representación}

En la aplicación intervienen tres actores:
\begin{itemize}
    \item \textbf{Usuario}. Utiliza las funciones de la aplicación. La necesidad
    de tener una cuenta o una sesión iniciada se indica en cada caso.
    \item \textbf{Servicio de correo}. Envía los enlaces para verificar la
    cuenta y restablecer la contraseña.
    \item \textbf{Catálogo externo de juegos}. Proporciona información para
    completar las búsquedas de juegos.
\end{itemize}

La figura~\ref{fig:cu:general} recoge los casos agrupados por función.
El diagrama general de la figura~\ref{fig:cu:global} muestra las relaciones
entre ellos con el \textbf{lenguaje unificado de modelado}
(\acs{UML}, del inglés \textit{Unified Modeling Language})\acused{UML}~\cite{omg_uml_251}.

Las conexiones con los actores se dibujan con líneas continuas. Sus flechas
indican el sentido de la solicitud: del usuario al caso de uso, o del caso
al servicio externo. Esta convención facilita la lectura del diagrama.

\newcommand{\refCasoUso}[1]{\hyperref[tab:cu:#1]{#1}}
\newlist{pasosCU}{enumerate}{1}
\setlist[pasosCU]{label=\arabic*.,leftmargin=1.6em,labelsep=0.5em,nosep,itemsep=2pt}
\newlist{alternativasCU}{itemize}{1}
\setlist[alternativasCU]{label=\textbullet,leftmargin=1.6em,labelsep=0.5em,nosep,itemsep=2pt}
\newfixedcaption{\captionFiguraCU}{figure}

\begin{figure}[p]
    \centering
    \includegraphics[width=\linewidth,height=0.88\textheight,keepaspectratio]{anexos/especificacion-de-requisitos/uml-general.pdf}
    \caption{Resumen de los casos de uso de \textit{Manualito}.}
    \label{fig:cu:general}
\end{figure}
\clearpage

\includepdf[
    pages=1,
    scale=0.94,
    offset=0 8,
    pagecommand={\thispagestyle{empty}\AddToShipoutPictureFG*{\AtPageLowerLeft{
        \put(0,20){\makebox[\paperwidth]{\parbox{0.9\paperwidth}{
            \captionFiguraCU{Diagrama general de casos de uso de \textit{Manualito}.}
            \label{fig:cu:global}
        }}}
    }}}
]{img/anexos/especificacion-de-requisitos/uml-global.pdf}

\subsection{Descripción de los casos de uso}

Cada uno de los casos de uso indica el autor, los actores, el objetivo, el inicio del caso,
las precondiciones, las acciones y el resultado esperado. Las alternativas
describen las decisiones y los errores que modifican la secuencia principal.

Esta organización permite comprobar el comportamiento y seguir su relación
con los requisitos, de acuerdo con la adaptación de
\acs{ISO}/\acs{IEC}/\acs{IEEE} 29148:2018 utilizada en el
proyecto~\cite{iso_29148_2018}.

\begingroup
\footnotesize
\setlength{\parskip}{0pt}
\SetTblrInner[longtblr]{measure=vbox,stretch=-1}
\SetTblrTemplate{conthead,contfoot,middlehead,lasthead}{empty}

\subsubsection{Cuenta y acceso}

% CU-S1
\begin{longtblr}[
    caption={\textbf{CU-S1}. Gestionar la cuenta.},
    label={tab:cu:CU-S1}
]{
    width=\linewidth,
    colspec={Q[l,wd=0.24\linewidth] X[l]},
    rows={valign=t},
    row{1}={bg=gray!18,font=\bfseries},
    column{1}={font=\bfseries},
    hline{1,Z}={0.8pt}, hline{2}={0.5pt},
    hborder{2-9,Y}={pagebreak=no},
    rowhead=1, leftsep=4pt, rightsep=4pt, rowsep=2.8pt
}
        CU-S1 & Gestionar la cuenta \\
        Autor & Ibai Moya Aroz \\
        Requisitos asociados & \refRequisito{RF-1.3}, \refRequisito{RF-1.4}, \refRequisito{RF-1.6}, \refRequisito{RF-1.7}, \refRequisito{RF-1.8}, \refRequisito{RF-1.9}, \refRequisito{RF-5.1}, \refRequisito{RF-5.2} \\
        Actores & Usuario. \\
        Descripción & El usuario consulta los datos y las opciones de su cuenta y decide si necesita realizar alguna operación. \\
        Inicio & El usuario abre el perfil, la sección de seguridad o los ajustes de su cuenta. \\
        Precondiciones & El usuario ha iniciado sesión. \\
        Acciones & \begin{pasosCU}
            \item El sistema muestra la información y las operaciones disponibles en la sección elegida.
            \item El usuario consulta los datos y decide si quiere realizar alguna de las operaciones disponibles.
            \item Si realiza una operación, el sistema muestra su resultado. El usuario puede continuar la gestión o terminarla.
        \end{pasosCU} \\
        Postcondiciones & El usuario ha consultado la información de su cuenta y los cambios aceptados quedan aplicados. Si cierra sesión o elimina la cuenta, la gestión termina mostrando el acceso público. \\
        Alternativas & \begin{alternativasCU}
            \item El usuario puede consultar la información y salir sin modificarla.
            \item Si cancela una operación o el sistema la rechaza, puede continuar la gestión sin dar por aplicado el cambio.
        \end{alternativasCU} \\
\end{longtblr}

\clearpage

% CU-01
\begin{longtblr}[
    caption={\textbf{CU-01}. Registrar una cuenta.},
    label={tab:cu:CU-01}
]{
    width=\linewidth,
    colspec={Q[l,wd=0.24\linewidth] X[l]},
    rows={valign=t},
    row{1}={bg=gray!18,font=\bfseries},
    column{1}={font=\bfseries},
    hline{1,Z}={0.8pt}, hline{2}={0.5pt},
    hborder{2-9,Y}={pagebreak=no},
    rowhead=1, leftsep=4pt, rightsep=4pt, rowsep=2.8pt
}
        CU-01 & Registrar una cuenta \\
        Autor & Ibai Moya Aroz \\
        Requisitos asociados & \refRequisito{RF-1.1} \\
        Actores & Usuario. \\
        Descripción & El usuario crea una cuenta para acceder a las funciones de \textit{Manualito} que requieren identificación. \\
        Inicio & El usuario selecciona la opción de crear una cuenta. \\
        Precondiciones & El usuario no ha iniciado sesión y dispone de un correo electrónico al que puede acceder. \\
        Acciones & \begin{pasosCU}
            \item El usuario introduce su nombre de usuario, correo electrónico y contraseña, y confirma esta última.
            \item Acepta la política de privacidad y solicita crear la cuenta.
            \item El sistema valida los datos, crea la cuenta e inicia la sesión.
            \item Mediante \refCasoUso{CU-31}, el sistema encarga el envío del enlace de verificación. El registro termina y el usuario puede acceder a la aplicación.
        \end{pasosCU} \\
        Postcondiciones & La cuenta queda creada y la sesión iniciada, con la verificación pendiente. El usuario puede verificarla después mediante \refCasoUso{CU-04}. \\
        Alternativas & \begin{alternativasCU}
            \item Si los datos no son válidos, el nombre o el correo ya están registrados, o falta la aceptación de la política, no se completa el registro.
            \item Si falla la entrega del correo, la cuenta sigue pendiente de verificación. El usuario puede solicitar otro envío mediante \refCasoUso{CU-31}.
        \end{alternativasCU} \\
\end{longtblr}

\clearpage

% CU-02
\begin{longtblr}[
    caption={\textbf{CU-02}. Iniciar sesión.},
    label={tab:cu:CU-02}
]{
    width=\linewidth,
    colspec={Q[l,wd=0.24\linewidth] X[l]},
    rows={valign=t},
    row{1}={bg=gray!18,font=\bfseries},
    column{1}={font=\bfseries},
    hline{1,Z}={0.8pt}, hline{2}={0.5pt},
    hborder{2-9,Y}={pagebreak=no},
    rowhead=1, leftsep=4pt, rightsep=4pt, rowsep=2.8pt
}
        CU-02 & Iniciar sesión \\
        Autor & Ibai Moya Aroz \\
        Requisitos asociados & \refRequisito{RF-1.2} \\
        Actores & Usuario. \\
        Descripción & El usuario accede a su cuenta mediante sus credenciales. \\
        Inicio & El usuario abre el formulario de acceso a su cuenta. \\
        Precondiciones & El usuario tiene una cuenta activa y no ha iniciado sesión. \\
        Acciones & \begin{pasosCU}
            \item El usuario introduce su correo electrónico o nombre de usuario y su contraseña.
            \item Solicita iniciar sesión y el sistema comprueba las credenciales.
            \item Si son válidas, el sistema inicia la sesión y permite acceder a las funciones de la cuenta.
        \end{pasosCU} \\
        Postcondiciones & El usuario queda identificado y puede utilizar las funciones que requieren una sesión iniciada. \\
        Alternativas & \begin{alternativasCU}
            \item Si las credenciales no son válidas, el sistema rechaza el acceso y el usuario puede corregirlas o abandonar el intento.
            \item La falta de verificación de la cuenta no impide iniciar sesión.
        \end{alternativasCU} \\
\end{longtblr}

\clearpage

% CU-03
\begin{longtblr}[
    caption={\textbf{CU-03}. Cerrar sesión.},
    label={tab:cu:CU-03}
]{
    width=\linewidth,
    colspec={Q[l,wd=0.24\linewidth] X[l]},
    rows={valign=t},
    row{1}={bg=gray!18,font=\bfseries},
    column{1}={font=\bfseries},
    hline{1,Z}={0.8pt}, hline{2}={0.5pt},
    hborder{2-9,Y}={pagebreak=no},
    rowhead=1, leftsep=4pt, rightsep=4pt, rowsep=2.8pt
}
        CU-03 & Cerrar sesión \\
        Autor & Ibai Moya Aroz \\
        Requisitos asociados & \refRequisito{RF-1.3} \\
        Actores & Usuario. \\
        Descripción & El usuario cierra la sesión que está utilizando. \\
        Inicio & El usuario solicita cerrar la sesión desde el perfil o los ajustes. \\
        Precondiciones & El usuario ha iniciado sesión. \\
        Acciones & \begin{pasosCU}
            \item El sistema revoca la sesión.
            \item El sistema muestra el acceso público.
        \end{pasosCU} \\
        Postcondiciones & La sesión cerrada deja de permitir el acceso a las funciones protegidas. Para volver a utilizarlas, el usuario debe iniciar sesión. \\
        Alternativas & Si la sesión ya ha caducado, el usuario permanece en el acceso público. \\
\end{longtblr}

\clearpage

% CU-04
\begin{longtblr}[
    caption={\textbf{CU-04}. Verificar la cuenta.},
    label={tab:cu:CU-04}
]{
    width=\linewidth,
    colspec={Q[l,wd=0.24\linewidth] X[l]},
    rows={valign=t},
    row{1}={bg=gray!18,font=\bfseries},
    column{1}={font=\bfseries},
    hline{1,Z}={0.8pt}, hline{2}={0.5pt},
    hborder{2-9,Y}={pagebreak=no},
    rowhead=1, leftsep=4pt, rightsep=4pt, rowsep=2.8pt
}
        CU-04 & Verificar la cuenta \\
        Autor & Ibai Moya Aroz \\
        Requisitos asociados & \refRequisito{RF-1.4} \\
        Actores & Usuario. \\
        Descripción & El usuario confirma que tiene acceso al correo electrónico asociado a su cuenta. \\
        Inicio & El usuario abre el enlace de verificación recibido por correo. \\
        Precondiciones & La cuenta está activa y el usuario ha recibido un enlace de verificación. No necesita haber iniciado sesión. \\
        Acciones & \begin{pasosCU}
            \item El sistema comprueba que el enlace permite verificar la cuenta.
            \item Marca la cuenta como verificada e invalida el enlace utilizado.
            \item El sistema muestra el resultado de la verificación.
        \end{pasosCU} \\
        Postcondiciones & La cuenta queda verificada y el enlace no permite una nueva verificación. \\
        Alternativas & \begin{alternativasCU}
            \item El sistema rechaza los enlaces que no reconoce. Si la cuenta sigue pendiente de verificación, también rechaza los enlaces caducados o utilizados.
            \item Si la cuenta ya está verificada y el sistema reconoce el enlace, conserva ese estado.
            \item Si necesita otro enlace, el usuario puede solicitarlo mediante \refCasoUso{CU-31}.
        \end{alternativasCU} \\
\end{longtblr}

\clearpage

% CU-05
\begin{longtblr}[
    caption={\textbf{CU-05}. Restablecer la contraseña.},
    label={tab:cu:CU-05}
]{
    width=\linewidth,
    colspec={Q[l,wd=0.24\linewidth] X[l]},
    rows={valign=t},
    row{1}={bg=gray!18,font=\bfseries},
    column{1}={font=\bfseries},
    hline{1,Z}={0.8pt}, hline{2}={0.5pt},
    hborder{2-9,Y}={pagebreak=no},
    rowhead=1, leftsep=4pt, rightsep=4pt, rowsep=2.8pt
}
        CU-05 & Restablecer la contraseña \\
        Autor & Ibai Moya Aroz \\
        Requisitos asociados & \refRequisito{RF-1.5} \\
        Actores & Usuario, Servicio de correo. \\
        Descripción & El usuario recupera el acceso a su cuenta cuando no recuerda la contraseña. \\
        Inicio & El usuario selecciona la opción de restablecer una contraseña olvidada. \\
        Precondiciones & La cuenta está activa y el usuario tiene acceso a su correo electrónico. \\
        Acciones & \begin{pasosCU}
            \item El usuario solicita restablecer la contraseña e indica el correo electrónico de su cuenta.
            \item El sistema genera un enlace de un solo uso y encarga su envío al servicio de correo.
            \item El usuario abre el enlace recibido e introduce y confirma una nueva contraseña.
            \item El sistema valida el enlace y la contraseña, guarda el cambio, invalida el enlace y revoca las sesiones anteriores.
            \item El sistema confirma el cambio y permite volver al formulario de acceso.
        \end{pasosCU} \\
        Postcondiciones & La contraseña queda actualizada y el correo queda verificado si todavía no lo estaba. El enlace utilizado y las sesiones anteriores dejan de ser válidos. \\
        Alternativas & \begin{alternativasCU}
            \item La respuesta a la solicitud no revela si el correo electrónico está registrado.
            \item Si el enlace no es válido, ha caducado o ya se ha utilizado, el sistema rechaza el restablecimiento.
            \item Si la nueva contraseña no cumple las condiciones o su confirmación no coincide, el usuario debe corregirla antes de continuar.
        \end{alternativasCU} \\
\end{longtblr}


\clearpage

% CU-06
\begin{longtblr}[
    caption={\textbf{CU-06}. Editar el perfil.},
    label={tab:cu:CU-06}
]{
    width=\linewidth,
    colspec={Q[l,wd=0.24\linewidth] X[l]},
    rows={valign=t},
    row{1}={bg=gray!18,font=\bfseries},
    column{1}={font=\bfseries},
    hline{1,Z}={0.8pt}, hline{2}={0.5pt},
    hborder{2-9,Y}={pagebreak=no},
    rowhead=1, leftsep=4pt, rightsep=4pt, rowsep=2.8pt
}
        CU-06 & Editar el perfil \\
        Autor & Ibai Moya Aroz \\
        Requisitos asociados & \refRequisito{RF-1.6} \\
        Actores & Usuario. \\
        Descripción & El usuario actualiza los datos con los que se identifica en su cuenta. \\
        Inicio & El usuario abre la edición del perfil desde su cuenta. \\
        Precondiciones & El usuario ha iniciado sesión. \\
        Acciones & \begin{pasosCU}
            \item El usuario modifica el nombre de usuario, el correo electrónico o la imagen de perfil.
            \item Solicita guardar los cambios y el sistema comprueba los datos modificados.
            \item El sistema guarda los cambios aceptados y muestra el perfil actualizado.
        \end{pasosCU} \\
        Postcondiciones & El perfil muestra los cambios guardados. Si se ha modificado el correo electrónico, la cuenta queda pendiente de verificarlo. \\
        Alternativas & \begin{alternativasCU}
            \item El usuario puede cancelar la edición sin guardar los cambios.
            \item Si algún dato no es válido o el nombre o el correo ya están registrados, el sistema rechaza el cambio y permite corregirlo.
        \end{alternativasCU} \\
\end{longtblr}

\clearpage

% CU-07
\begin{longtblr}[
    caption={\textbf{CU-07}. Cambiar la contraseña.},
    label={tab:cu:CU-07}
]{
    width=\linewidth,
    colspec={Q[l,wd=0.24\linewidth] X[l]},
    rows={valign=t},
    row{1}={bg=gray!18,font=\bfseries},
    column{1}={font=\bfseries},
    hline{1,Z}={0.8pt}, hline{2}={0.5pt},
    hborder{2-9,Y}={pagebreak=no},
    rowhead=1, leftsep=4pt, rightsep=4pt, rowsep=2.8pt
}
        CU-07 & Cambiar la contraseña \\
        Autor & Ibai Moya Aroz \\
        Requisitos asociados & \refRequisito{RF-1.7} \\
        Actores & Usuario. \\
        Descripción & El usuario sustituye la contraseña de su cuenta y conserva la sesión desde la que realiza el cambio. \\
        Inicio & El usuario accede a la sección de seguridad para cambiar la contraseña. \\
        Precondiciones & El usuario ha iniciado sesión y conoce su contraseña actual. \\
        Acciones & \begin{pasosCU}
            \item El usuario introduce la contraseña actual y la nueva, que debe confirmar.
            \item Solicita el cambio y el sistema comprueba la contraseña actual y valida la nueva.
            \item El sistema guarda la nueva contraseña, revoca las demás sesiones y confirma el cambio.
        \end{pasosCU} \\
        Postcondiciones & La contraseña queda actualizada y solo permanece activa la sesión desde la que se ha realizado el cambio. \\
        Alternativas & Si la contraseña actual es incorrecta, la nueva no cumple las condiciones o su confirmación no coincide, el sistema impide el cambio. \\
\end{longtblr}

\clearpage

% CU-08
\begin{longtblr}[
    caption={\textbf{CU-08}. Eliminar la cuenta.},
    label={tab:cu:CU-08}
]{
    width=\linewidth,
    colspec={Q[l,wd=0.24\linewidth] X[l]},
    rows={valign=t},
    row{1}={bg=gray!18,font=\bfseries},
    column{1}={font=\bfseries},
    hline{1,Z}={0.8pt}, hline{2}={0.5pt},
    hborder{2-9,Y}={pagebreak=no},
    rowhead=1, leftsep=4pt, rightsep=4pt, rowsep=2.8pt
}
        CU-08 & Eliminar la cuenta \\
        Autor & Ibai Moya Aroz \\
        Requisitos asociados & \refRequisito{RF-1.8} \\
        Actores & Usuario. \\
        Descripción & El usuario elimina su cuenta y retira el acceso a la información asociada a ella. \\
        Inicio & El usuario solicita eliminar su cuenta desde los ajustes o la sección de seguridad. \\
        Precondiciones & El usuario ha iniciado sesión en la cuenta que desea eliminar. \\
        Acciones & \begin{pasosCU}
            \item El usuario revisa la información que dejará de estar disponible e introduce su nombre de usuario para confirmar el borrado.
            \item El sistema comprueba la confirmación, revoca las sesiones y hace inaccesibles la cuenta y sus recursos.
            \item El sistema inicia la limpieza de los recursos asociados y muestra el acceso público.
        \end{pasosCU} \\
        Postcondiciones & La cuenta, los manuales, las conversaciones y las valoraciones asociadas dejan de ser accesibles, y sus sesiones dejan de ser válidas. \\
        Alternativas & \begin{alternativasCU}
            \item El usuario puede cancelar antes de confirmar el borrado y continuar utilizando la cuenta.
            \item Si el nombre introducido no coincide con el de la cuenta, el sistema impide la eliminación.
            \item Si la solicitud falla, el sistema informa del error sin confirmar que la cuenta se haya eliminado.
        \end{alternativasCU} \\
\end{longtblr}

\clearpage

% CU-09
\begin{longtblr}[
    caption={\textbf{CU-09}. Consultar la actividad personal.},
    label={tab:cu:CU-09}
]{
    width=\linewidth,
    colspec={Q[l,wd=0.24\linewidth] X[l]},
    rows={valign=t},
    row{1}={bg=gray!18,font=\bfseries},
    column{1}={font=\bfseries},
    hline{1,Z}={0.8pt}, hline{2}={0.5pt},
    hborder{2-9,Y}={pagebreak=no},
    rowhead=1, leftsep=4pt, rightsep=4pt, rowsep=2.8pt
}
        CU-09 & Consultar la actividad personal \\
        Autor & Ibai Moya Aroz \\
        Requisitos asociados & \refRequisito{RF-1.9} \\
        Actores & Usuario. \\
        Descripción & El usuario consulta un resumen de la actividad que ha realizado con su cuenta. \\
        Inicio & El usuario abre su perfil para consultar la actividad personal. \\
        Precondiciones & El usuario ha iniciado sesión. \\
        Acciones & \begin{pasosCU}
            \item El sistema obtiene los recuentos de actividad de la cuenta.
            \item El sistema muestra, junto a los datos del perfil, el número de juegos con actividad, manuales subidos y conversaciones creadas.
        \end{pasosCU} \\
        Postcondiciones & El usuario ha consultado los recuentos de su actividad personal. \\
        Alternativas & Si no hay actividad de alguno de estos tipos, su recuento es cero. \\
\end{longtblr}

\clearpage

% CU-31
\begin{longtblr}[
    caption={\textbf{CU-31}. Enviar el enlace de verificación.},
    label={tab:cu:CU-31}
]{
    width=\linewidth,
    colspec={Q[l,wd=0.24\linewidth] X[l]},
    rows={valign=t},
    row{1}={bg=gray!18,font=\bfseries},
    column{1}={font=\bfseries},
    hline{1,Z}={0.8pt}, hline{2}={0.5pt},
    hborder{2-9,Y}={pagebreak=no},
    rowhead=1, leftsep=4pt, rightsep=4pt, rowsep=2.8pt
}
        CU-31 & Enviar el enlace de verificación \\
        Autor & Ibai Moya Aroz \\
        Requisitos asociados & \refRequisito{RF-1.4} \\
        Actores & Usuario, Servicio de correo. \\
        Descripción & El sistema proporciona un enlace para que el usuario pueda verificar su cuenta. \\
        Inicio & Se crea una cuenta, se cambia su correo o el usuario solicita otro enlace de verificación. \\
        Precondiciones & La cuenta está activa y su correo está pendiente de verificación. \\
        Acciones & \begin{pasosCU}
            \item El sistema obtiene el correo de destino y comprueba que procede emitir un enlace.
            \item Genera un enlace de verificación vigente e invalida los anteriores cuando los haya.
            \item Encarga el envío al servicio de correo y continúa la operación que lo solicitó o responde a la petición de reenvío.
        \end{pasosCU} \\
        Postcondiciones & Se emite un enlace y se encarga su envío. La cuenta sigue pendiente de verificación hasta que el usuario utilice un enlace válido en \refCasoUso{CU-04}. \\
        Alternativas & \begin{alternativasCU}
            \item La respuesta a una solicitud pública no revela si existe una cuenta con ese correo.
            \item Si la cuenta ya está verificada, no está activa o no ha transcurrido la espera entre solicitudes, no se emite otro enlace.
            \item Si falla la entrega, la cuenta conserva la verificación pendiente.
        \end{alternativasCU} \\
\end{longtblr}

\clearpage

\subsubsection{Juegos y biblioteca}

% CU-10
\begin{longtblr}[
    caption={\textbf{CU-10}. Buscar juegos.},
    label={tab:cu:CU-10}
]{
    width=\linewidth,
    colspec={Q[l,wd=0.24\linewidth] X[l]},
    rows={valign=t},
    row{1}={bg=gray!18,font=\bfseries},
    column{1}={font=\bfseries},
    hline{1,Z}={0.8pt}, hline{2}={0.5pt},
    hborder{2-9,Y}={pagebreak=no},
    rowhead=1, leftsep=4pt, rightsep=4pt, rowsep=2.8pt
}
        CU-10 & Buscar juegos \\
        Autor & Ibai Moya Aroz \\
        Requisitos asociados & \refRequisito{RF-2.1} \\
        Actores & Usuario, Catálogo externo de juegos. \\
        Descripción & Permite localizar un juego por su nombre para conocerlo o asociarle un manual. \\
        Inicio & El usuario abre la búsqueda durante la exploración o al añadir un manual. \\
        Precondiciones & El usuario ha iniciado sesión. \\
        Acciones & \begin{pasosCU}
            \item El usuario introduce el nombre del juego que busca.
            \item El sistema busca coincidencias y muestra los resultados.
            \item El usuario puede seleccionar un resultado para consultar el juego durante la exploración o para continuar añadiendo un manual.
        \end{pasosCU} \\
        Postcondiciones & El sistema muestra los resultados de la búsqueda o indica que no hay coincidencias. Si selecciona un juego durante una subida, ese juego queda elegido para asociarle el manual. \\
        Alternativas & \begin{alternativasCU}
            \item Si no hay coincidencias, el sistema lo indica y el usuario puede cambiar la búsqueda o terminarla.
            \item Durante la subida de un manual, una búsqueda sin coincidencias también permite añadir el juego manualmente.
        \end{alternativasCU} \\
\end{longtblr}

\clearpage

% CU-11
\begin{longtblr}[
    caption={\textbf{CU-11}. Añadir un juego manualmente.},
    label={tab:cu:CU-11}
]{
    width=\linewidth,
    colspec={Q[l,wd=0.24\linewidth] X[l]},
    rows={valign=t},
    row{1}={bg=gray!18,font=\bfseries},
    column{1}={font=\bfseries},
    hline{1,Z}={0.8pt}, hline{2}={0.5pt},
    hborder{2-9,Y}={pagebreak=no},
    rowhead=1, leftsep=4pt, rightsep=4pt, rowsep=2.8pt
}
        CU-11 & Añadir un juego manualmente \\
        Autor & Ibai Moya Aroz \\
        Requisitos asociados & \refRequisito{RF-2.2} \\
        Actores & Usuario. \\
        Descripción & Permite registrar un juego que no aparece en la búsqueda para poder asociarle el manual que se está añadiendo. \\
        Inicio & El usuario elige añadir el juego tras una búsqueda sin coincidencias. \\
        Precondiciones & El usuario ha iniciado sesión y la búsqueda de \refCasoUso{CU-10} no ha devuelto coincidencias durante \refCasoUso{CU-16}. \\
        Acciones & \begin{pasosCU}
            \item El usuario confirma el nombre del juego que quiere añadir.
            \item El sistema valida el nombre, registra el juego y lo selecciona para continuar la subida.
        \end{pasosCU} \\
        Postcondiciones & El juego queda registrado y seleccionado para asociarle el manual. \\
        Alternativas & \begin{alternativasCU}
            \item Si el nombre no es válido, el usuario debe corregirlo antes de añadir el juego.
            \item El usuario puede cancelar la creación y volver a buscar un juego existente.
        \end{alternativasCU} \\
\end{longtblr}

\clearpage

% CU-12
\begin{longtblr}[
    caption={\textbf{CU-12}. Seguir o dejar de seguir un juego.},
    label={tab:cu:CU-12}
]{
    width=\linewidth,
    colspec={Q[l,wd=0.24\linewidth] X[l]},
    rows={valign=t},
    row{1}={bg=gray!18,font=\bfseries},
    column{1}={font=\bfseries},
    hline{1,Z}={0.8pt}, hline{2}={0.5pt},
    hborder{2-9,Y}={pagebreak=no},
    rowhead=1, leftsep=4pt, rightsep=4pt, rowsep=2.8pt
}
        CU-12 & Seguir o dejar de seguir un juego \\
        Autor & Ibai Moya Aroz \\
        Requisitos asociados & \refRequisito{RF-2.3}, \refRequisito{RF-2.3.1} \\
        Actores & Usuario. \\
        Descripción & Permite decidir si un juego debe aparecer entre los juegos seguidos de la biblioteca personal. \\
        Inicio & El usuario decide seguir el juego que está consultando o dejar de seguirlo. \\
        Precondiciones & El usuario ha iniciado sesión y el juego existe. \\
        Acciones & \begin{pasosCU}
            \item Durante la consulta del perfil, el usuario elige seguir el juego o dejar de seguirlo.
            \item El sistema guarda la decisión, actualiza la biblioteca personal e indica si el usuario sigue el juego.
        \end{pasosCU} \\
        Postcondiciones & La decisión queda guardada y el perfil indica si el usuario sigue el juego. \\
        Alternativas & Si el usuario decide dejar de seguir el juego, una actividad posterior no vuelve a seguirlo automáticamente. \\
\end{longtblr}

\clearpage

% CU-13
\begin{longtblr}[
    caption={\textbf{CU-13}. Consultar la biblioteca personal.},
    label={tab:cu:CU-13}
]{
    width=\linewidth,
    colspec={Q[l,wd=0.24\linewidth] X[l]},
    rows={valign=t},
    row{1}={bg=gray!18,font=\bfseries},
    column{1}={font=\bfseries},
    hline{1,Z}={0.8pt}, hline{2}={0.5pt},
    hborder{2-9,Y}={pagebreak=no},
    rowhead=1, leftsep=4pt, rightsep=4pt, rowsep=2.8pt
}
        CU-13 & Consultar la biblioteca personal \\
        Autor & Ibai Moya Aroz \\
        Requisitos asociados & \refRequisito{RF-2.4}, \refRequisito{RF-2.4.1}, \refRequisito{RF-2.4.2} \\
        Actores & Usuario. \\
        Descripción & Permite localizar los juegos seguidos y los manuales propios y consultar la información de cada uno. \\
        Inicio & El usuario accede a la biblioteca personal. \\
        Precondiciones & El usuario ha iniciado sesión. \\
        Acciones & \begin{pasosCU}
            \item El usuario elige consultar los juegos o los manuales.
            \item El sistema muestra los juegos seguidos o los manuales propios.
            \item El usuario puede filtrar los manuales y el sistema muestra las coincidencias.
            \item El usuario puede seleccionar una entrada para consultar los detalles del juego, el contenido del manual o su estado de procesamiento.
        \end{pasosCU} \\
        Postcondiciones & La biblioteca muestra los juegos seguidos o los manuales propios según la selección y los filtros aplicados. Si el usuario elimina un manual, la lista deja de mostrarlo. \\
        Alternativas & \begin{alternativasCU}
            \item Si la lista está vacía o el filtro no devuelve resultados, el sistema lo indica.
            \item El usuario puede eliminar un manual propio durante la consulta. Tras confirmar la eliminación, la biblioteca muestra los manuales restantes o indica que no hay ninguno.
        \end{alternativasCU} \\
\end{longtblr}

\clearpage

% CU-14
\begin{longtblr}[
    caption={\textbf{CU-14}. Consultar el perfil de un juego.},
    label={tab:cu:CU-14}
]{
    width=\linewidth,
    colspec={Q[l,wd=0.24\linewidth] X[l]},
    rows={valign=t},
    row{1}={bg=gray!18,font=\bfseries},
    column{1}={font=\bfseries},
    hline{1,Z}={0.8pt}, hline{2}={0.5pt},
    hborder{2-9,Y}={pagebreak=no},
    rowhead=1, leftsep=4pt, rightsep=4pt, rowsep=2.8pt
}
        CU-14 & Consultar el perfil de un juego \\
        Autor & Ibai Moya Aroz \\
        Requisitos asociados & \refRequisito{RF-2.5} \\
        Actores & Usuario. \\
        Descripción & Permite conocer un juego, revisar la actividad personal asociada y ampliar la consulta mediante sus manuales y conversaciones. \\
        Inicio & El usuario selecciona un juego para consultar su perfil. \\
        Precondiciones & El usuario ha iniciado sesión y el juego existe. \\
        Acciones & \begin{pasosCU}
            \item El sistema muestra los datos del juego, los manuales accesibles, la valoración propia y el número de conversaciones del usuario, e indica si sigue el juego.
            \item El usuario consulta esta información y puede profundizar en las reglas, los manuales o las conversaciones antes de terminar.
        \end{pasosCU} \\
        Postcondiciones & El perfil refleja la información accesible y los cambios realizados durante la consulta. \\
        Alternativas & \begin{alternativasCU}
            \item Sin manuales utilizables no se generan explicaciones ni respuestas nuevas, pero el perfil sigue siendo consultable.
            \item Los manuales compartidos pueden aportar contenido para consultas y abrirse en el visor en modo de solo lectura.
        \end{alternativasCU} \\
\end{longtblr}

\clearpage

% CU-15
\begin{longtblr}[
    caption={\textbf{CU-15}. Gestionar la valoración de un juego.},
    label={tab:cu:CU-15}
]{
    width=\linewidth,
    colspec={Q[l,wd=0.24\linewidth] X[l]},
    rows={valign=t},
    row{1}={bg=gray!18,font=\bfseries},
    column{1}={font=\bfseries},
    hline{1,Z}={0.8pt}, hline{2}={0.5pt},
    hborder{2-9,Y}={pagebreak=no},
    rowhead=1, leftsep=4pt, rightsep=4pt, rowsep=2.8pt
}
        CU-15 & Gestionar la valoración de un juego \\
        Autor & Ibai Moya Aroz \\
        Requisitos asociados & \refRequisito{RF-2.6}, \refRequisito{RF-2.3.2} \\
        Actores & Usuario. \\
        Descripción & Permite conservar una valoración personal del juego mediante una puntuación y una nota opcional. \\
        Inicio & El usuario abre la opción de valoración del perfil del juego. \\
        Precondiciones & El usuario ha iniciado sesión y el juego existe. \\
        Acciones & \begin{pasosCU}
            \item Desde el perfil del juego, el usuario selecciona una puntuación de uno a cinco y puede añadir una nota.
            \item El usuario solicita guardar la valoración y el sistema la registra o actualiza.
            \item El perfil muestra la valoración guardada.
        \end{pasosCU} \\
        Postcondiciones & La valoración queda creada, modificada o eliminada y el perfil refleja el resultado. \\
        Alternativas & \begin{alternativasCU}
            \item Si ya existe una valoración, el usuario puede modificarla o confirmar su eliminación.
            \item El usuario puede cancelar sin guardar los cambios.
            \item Si la valoración es la primera actividad sobre el juego, el sistema lo sigue automáticamente cuando no existe una decisión explícita anterior.
        \end{alternativasCU} \\
\end{longtblr}

\clearpage

\subsubsection{Manuales}

% CU-16
\begin{longtblr}[
    caption={\textbf{CU-16}. Añadir un manual.},
    label={tab:cu:CU-16}
]{
    width=\linewidth,
    colspec={Q[l,wd=0.24\linewidth] X[l]},
    rows={valign=t},
    row{1}={bg=gray!18,font=\bfseries},
    column{1}={font=\bfseries},
    hline{1,Z}={0.8pt}, hline{2}={0.5pt},
    hborder{2-9,Y}={pagebreak=no},
    rowhead=1, leftsep=4pt, rightsep=4pt, rowsep=2.8pt
}
        CU-16 & Añadir un manual \\
        Autor & Ibai Moya Aroz \\
        Requisitos asociados & \refRequisito{RF-3.1}, \refRequisito{RF-3.2}, \refRequisito{RF-3.3}, \refRequisito{RF-3.10}, \refRequisito{RF-2.3.2} \\
        Actores & Usuario. \\
        Descripción & Permite asociar un manual a un juego y preparar su contenido para utilizarlo en las consultas. \\
        Inicio & El usuario solicita añadir un manual, desde el perfil de un juego o desde la opción de subida. \\
        Precondiciones & El usuario ha iniciado sesión y dispone del manual en un documento \acs{PDF} o en imágenes. \\
        Acciones & \begin{pasosCU}
            \item El usuario acepta el juego preseleccionado o selecciona uno.
            \item El usuario añade el documento o las imágenes y elige si el manual será privado o compartido. Si añade varias imágenes, puede revisarlas, ordenarlas y retirar las que no sean útiles.
            \item Si comparte el manual, elige si quiere mostrar su nombre. Por defecto aparece como \guillemotleft{}Anónimo\guillemotright{}.
            \item El usuario confirma la subida. El sistema valida los formatos, los tamaños y el número de páginas, y comprueba si el manual está duplicado.
            \item El sistema guarda la fuente aceptada y ejecuta \refCasoUso{CU-30} en segundo plano.
        \end{pasosCU} \\
        Postcondiciones & El manual aceptado queda asociado al juego con la visibilidad y la opción de mostrar el nombre elegidas, y el juego empieza a seguirse automáticamente cuando corresponde. Su contenido queda preparado para consultas si el procesamiento obtiene un resultado utilizable. En caso contrario, se registra el fallo o la ausencia de texto útil. \\
        Alternativas & \begin{alternativasCU}
            \item El usuario puede cancelar antes de enviar el manual.
            \item El sistema rechaza los archivos que incumplen los límites y los manuales duplicados antes de procesarlos.
            \item El usuario puede seguir utilizando la aplicación mientras se procesa el manual.
        \end{alternativasCU} \\
\end{longtblr}


\clearpage

% CU-17
\begin{longtblr}[
    caption={\textbf{CU-17}. Consultar el estado de un manual.},
    label={tab:cu:CU-17}
]{
    width=\linewidth,
    colspec={Q[l,wd=0.24\linewidth] X[l]},
    rows={valign=t},
    row{1}={bg=gray!18,font=\bfseries},
    column{1}={font=\bfseries},
    hline{1,Z}={0.8pt}, hline{2}={0.5pt},
    hborder{2-9,Y}={pagebreak=no},
    rowhead=1, leftsep=4pt, rightsep=4pt, rowsep=2.8pt
}
        CU-17 & Consultar el estado de un manual \\
        Autor & Ibai Moya Aroz \\
        Requisitos asociados & \refRequisito{RF-3.5} \\
        Actores & Usuario. \\
        Descripción & Permite conocer el avance del procesamiento de un manual y el resultado de cada página. \\
        Inicio & El usuario elige consultar el estado de uno de sus manuales. \\
        Precondiciones & El usuario ha iniciado sesión y el manual le pertenece. No es necesario que su procesamiento haya terminado. \\
        Acciones & \begin{pasosCU}
            \item El sistema muestra el estado guardado del manual y de sus páginas.
            \item El usuario revisa qué páginas siguen pendientes o en procesamiento y cuáles se han completado o han fallado.
        \end{pasosCU} \\
        Postcondiciones & El usuario conoce el estado guardado del manual y de sus páginas. \\
        Alternativas & \begin{alternativasCU}
            \item El usuario puede abandonar la consulta y revisar el estado más tarde.
            \item Los fallos y la ausencia de texto útil se muestran como estados distintos del contenido preparado para consultas.
        \end{alternativasCU} \\
\end{longtblr}

\clearpage

% CU-18
\begin{longtblr}[
    caption={\textbf{CU-18}. Consultar un manual.},
    label={tab:cu:CU-18}
]{
    width=\linewidth,
    colspec={Q[l,wd=0.24\linewidth] X[l]},
    rows={valign=t},
    row{1}={bg=gray!18,font=\bfseries},
    column{1}={font=\bfseries},
    hline{1,Z}={0.8pt}, hline{2}={0.5pt},
    hborder{2-9,Y}={pagebreak=no},
    rowhead=1, leftsep=4pt, rightsep=4pt, rowsep=2.8pt
}
        CU-18 & Consultar un manual \\
        Autor & Ibai Moya Aroz \\
        Requisitos asociados & \refRequisito{RF-3.6}, \refRequisito{RF-3.6.1}, \refRequisito{RF-3.6.2}, \refRequisito{RF-3.6.3}, \refRequisito{RF-3.10} \\
        Actores & Usuario. \\
        Descripción & Permite leer un manual propio o compartido y contrastar el texto de cada página con su imagen. La información de procesamiento solo se muestra al propietario. \\
        Inicio & El usuario selecciona un manual disponible o una referencia a una de sus páginas. \\
        Precondiciones & El usuario ha iniciado sesión. El manual no está eliminado y le pertenece o está compartido. \\
        Acciones & \begin{pasosCU}
            \item El sistema abre el manual en la página indicada. Si no se ha indicado ninguna o esa página no existe, abre la primera página disponible.
            \item El sistema muestra el texto extraído o corregido y la imagen disponible. Si el manual es propio, también muestra los datos de origen, confianza y reutilización de la página.
            \item El usuario puede navegar entre las páginas o buscar texto y desplazarse entre las coincidencias.
        \end{pasosCU} \\
        Postcondiciones & El usuario consulta las páginas elegidas. Si el manual pertenece a otra cuenta, no puede editarlo, reprocesarlo, eliminarlo ni cambiar sus opciones. \\
        Alternativas & \begin{alternativasCU}
            \item Si la búsqueda no devuelve coincidencias o falta información de procesamiento, el sistema lo indica.
            \item Si el manual está compartido, el propietario puede cambiar la opción de mostrar su nombre y confirmar la elección. El sistema la guarda y la aplica al manual y a las fuentes de las respuestas.
            \item Si cancela la confirmación o falla el guardado, se conserva la opción anterior.
        \end{alternativasCU} \\
\end{longtblr}

\clearpage

% CU-19
\begin{longtblr}[
    caption={\textbf{CU-19}. Corregir el texto de una página.},
    label={tab:cu:CU-19}
]{
    width=\linewidth,
    colspec={Q[l,wd=0.24\linewidth] X[l]},
    rows={valign=t},
    row{1}={bg=gray!18,font=\bfseries},
    column{1}={font=\bfseries},
    hline{1,Z}={0.8pt}, hline{2}={0.5pt},
    hborder{2-9,Y}={pagebreak=no},
    rowhead=1, leftsep=4pt, rightsep=4pt, rowsep=2.8pt
}
        CU-19 & Corregir el texto de una página \\
        Autor & Ibai Moya Aroz \\
        Requisitos asociados & \refRequisito{RF-3.7} \\
        Actores & Usuario. \\
        Descripción & Permite corregir errores o completar el texto de una página para utilizar la versión revisada en las consultas. \\
        Inicio & El usuario solicita editar el texto de la página que está consultando. \\
        Precondiciones & El usuario ha iniciado sesión, la página existe y el manual es propio, privado y no se está procesando. \\
        Acciones & \begin{pasosCU}
            \item El usuario modifica el texto de la página que está consultando y confirma el cambio.
            \item El sistema guarda la corrección y solicita actualizar el contenido indexado de esa página.
            \item El visor muestra el texto corregido para continuar la lectura.
        \end{pasosCU} \\
        Postcondiciones & El texto corregido queda guardado y se solicita su actualización para las consultas. El texto de las demás páginas se conserva. \\
        Alternativas & \begin{alternativasCU}
            \item El usuario puede cancelar sin guardar o añadir texto a una página que todavía no lo tenga.
            \item El sistema rechaza la edición de un manual ajeno, compartido o que se esté procesando.
            \item Si falla la solicitud de actualización del contenido indexado después de guardar, el texto corregido se conserva y el sistema informa del error.
        \end{alternativasCU} \\
\end{longtblr}


\clearpage

% CU-20
\begin{longtblr}[
    caption={\textbf{CU-20}. Reprocesar un manual.},
    label={tab:cu:CU-20}
]{
    width=\linewidth,
    colspec={Q[l,wd=0.24\linewidth] X[l]},
    rows={valign=t},
    row{1}={bg=gray!18,font=\bfseries},
    column{1}={font=\bfseries},
    hline{1,Z}={0.8pt}, hline{2}={0.5pt},
    hborder{2-9,Y}={pagebreak=no},
    rowhead=1, leftsep=4pt, rightsep=4pt, rowsep=2.8pt
}
        CU-20 & Reprocesar un manual \\
        Autor & Ibai Moya Aroz \\
        Requisitos asociados & \refRequisito{RF-3.8}, \refRequisito{RF-3.8.1}, \refRequisito{RF-3.8.2} \\
        Actores & Usuario. \\
        Descripción & Permite obtener de nuevo el contenido de un manual completo o de una página a partir de las fuentes guardadas. \\
        Inicio & El usuario solicita repetir el procesamiento desde el visor del manual. \\
        Precondiciones & El usuario ha iniciado sesión, el manual le pertenece, conserva sus fuentes y su estado permite repetir el procesamiento. \\
        Acciones & \begin{pasosCU}
            \item El usuario elige reprocesar el manual completo o una página y confirma la operación.
            \item El sistema ejecuta \refCasoUso{CU-30} en segundo plano sobre las páginas seleccionadas.
            \item El sistema guarda el nuevo resultado o el fallo y actualiza el estado del manual y de las páginas afectadas.
        \end{pasosCU} \\
        Postcondiciones & El nuevo resultado o el fallo queda registrado. Si se reprocesa una sola página, se conserva el resto del manual. \\
        Alternativas & \begin{alternativasCU}
            \item El usuario puede cancelar antes de confirmar la operación.
            \item Si el estado del manual impide reprocesarlo, el sistema rechaza la solicitud.
            \item El reprocesamiento puede sustituir las correcciones manuales y puede terminar con un error. No requiere que el procesamiento anterior haya fallado.
        \end{alternativasCU} \\
\end{longtblr}

\clearpage

% CU-21
\begin{longtblr}[
    caption={\textbf{CU-21}. Eliminar un manual.},
    label={tab:cu:CU-21}
]{
    width=\linewidth,
    colspec={Q[l,wd=0.24\linewidth] X[l]},
    rows={valign=t},
    row{1}={bg=gray!18,font=\bfseries},
    column{1}={font=\bfseries},
    hline{1,Z}={0.8pt}, hline{2}={0.5pt},
    hborder{2-9,Y}={pagebreak=no},
    rowhead=1, leftsep=4pt, rightsep=4pt, rowsep=2.8pt
}
        CU-21 & Eliminar un manual \\
        Autor & Ibai Moya Aroz \\
        Requisitos asociados & \refRequisito{RF-3.9} \\
        Actores & Usuario. \\
        Descripción & Permite retirar un manual propio y su contenido del acceso del usuario. \\
        Inicio & El usuario solicita eliminar un manual desde la biblioteca o desde su visor. \\
        Precondiciones & El usuario ha iniciado sesión. El manual le pertenece y no está eliminado. No es necesario que su procesamiento haya terminado. \\
        Acciones & \begin{pasosCU}
            \item El usuario confirma la eliminación del manual.
            \item El sistema retira el acceso al manual y sus páginas y solicita la limpieza de sus archivos y del contenido indexado.
            \item Si la eliminación se inició en la biblioteca, el sistema actualiza la lista. Si se inició en el visor, termina la consulta del manual y muestra la biblioteca.
        \end{pasosCU} \\
        Postcondiciones & El manual y sus páginas dejan de ser accesibles y se solicita la limpieza de sus recursos. El historial de conversaciones se conserva. \\
        Alternativas & \begin{alternativasCU}
            \item El usuario puede cancelar sin eliminar el manual.
            \item Si se elimina la última fuente utilizable de un juego, las conversaciones existentes se conservan y quedan en modo de solo lectura.
        \end{alternativasCU} \\
\end{longtblr}

\clearpage

% CU-30
\begin{longtblr}[
    caption={\textbf{CU-30}. Procesar el contenido del manual.},
    label={tab:cu:CU-30}
]{
    width=\linewidth,
    colspec={Q[l,wd=0.24\linewidth] X[l]},
    rows={valign=t},
    row{1}={bg=gray!18,font=\bfseries},
    column{1}={font=\bfseries},
    hline{1,Z}={0.8pt}, hline{2}={0.5pt},
    hborder{2-9,Y}={pagebreak=no},
    rowhead=1, leftsep=4pt, rightsep=4pt, rowsep=2.8pt
}
        CU-30 & Procesar el contenido del manual \\
        Autor & Ibai Moya Aroz \\
        Requisitos asociados & \refRequisito{RF-3.4}, \refRequisito{RF-3.4.1}, \refRequisito{RF-3.4.2}, \refRequisito{RF-3.4.3}, \refRequisito{RF-3.4.4}, \refRequisito{RF-3.4.5} \\
        Actores & Usuario. \\
        Descripción & Prepara el texto del manual para que pueda recuperarse al generar explicaciones y responder preguntas sobre el juego. \\
        Inicio & El sistema inicia el procesamiento previsto en \refCasoUso{CU-16} o en \refCasoUso{CU-20}. \\
        Precondiciones & Las fuentes están guardadas y las páginas que deben procesarse están identificadas. \\
        Acciones & \begin{pasosCU}
            \item El sistema utiliza el texto del documento cuando es aprovechable. En caso contrario, preprocesa las imágenes y extrae su texto mediante \ac{OCR}.
            \item El sistema aplica el posprocesado y guarda el texto resultante.
            \item El sistema divide el texto en fragmentos y los indexa para recuperarlos durante las consultas.
            \item El sistema actualiza el estado del manual y de sus páginas según el resultado del procesamiento.
        \end{pasosCU} \\
        Postcondiciones & El contenido utilizable queda preparado para las consultas. Cuando el procesamiento falla o no obtiene texto útil, el estado del manual y de sus páginas refleja esa situación. \\
        Alternativas & \begin{alternativasCU}
            \item Si encuentra una página idéntica ya procesada, el sistema reutiliza su resultado sin repetir la extracción ni el posprocesado.
            \item Si falla una etapa o no se obtiene texto útil, el sistema registra el estado correspondiente.
        \end{alternativasCU} \\
\end{longtblr}

\clearpage

\subsubsection{Consulta de reglas}

% CU-22
\begin{longtblr}[
    caption={\textbf{CU-22}. Consultar la explicación de un juego.},
    label={tab:cu:CU-22}
]{
    width=\linewidth,
    colspec={Q[l,wd=0.24\linewidth] X[l]},
    rows={valign=t},
    row{1}={bg=gray!18,font=\bfseries},
    column{1}={font=\bfseries},
    hline{1,Z}={0.8pt}, hline{2}={0.5pt},
    hborder{2-9,Y}={pagebreak=no},
    rowhead=1, leftsep=4pt, rightsep=4pt, rowsep=2.8pt
}
        CU-22 & Consultar la explicación de un juego \\
        Autor & Ibai Moya Aroz \\
        Requisitos asociados & \refRequisito{RF-4.1} \\
        Actores & Usuario. \\
        Descripción & Permite conocer la preparación, el desarrollo y las condiciones de victoria de un juego a partir de sus manuales. \\
        Inicio & El usuario abre el perfil de un juego con manuales utilizables. \\
        Precondiciones & El usuario ha iniciado sesión y el juego dispone de manuales que puede utilizar para las consultas. \\
        Acciones & \begin{pasosCU}
            \item El sistema solicita automáticamente la explicación del juego.
            \item El sistema recupera la explicación conservada si sus fuentes no han cambiado. Si no dispone de una explicación vigente, la genera en segundo plano.
            \item El usuario consulta los apartados de la explicación a medida que están disponibles.
        \end{pasosCU} \\
        Postcondiciones & La explicación generada puede leerse desde el perfil del juego y se conserva para consultas posteriores mientras no cambien sus fuentes. \\
        Alternativas & \begin{alternativasCU}
            \item Si no hay manuales utilizables, se informa de que faltan fuentes y no se genera la explicación.
            \item Si la generación falla antes de obtener contenido, el sistema informa del error. Si ya hay apartados disponibles, conserva su contenido. El resto del perfil sigue disponible.
        \end{alternativasCU} \\
\end{longtblr}

\clearpage

% CU-23
\begin{longtblr}[
    caption={\textbf{CU-23}. Consultar una duda sobre las reglas.},
    label={tab:cu:CU-23}
]{
    width=\linewidth,
    colspec={Q[l,wd=0.24\linewidth] X[l]},
    rows={valign=t},
    row{1}={bg=gray!18,font=\bfseries},
    column{1}={font=\bfseries},
    hline{1,Z}={0.8pt}, hline{2}={0.5pt},
    hborder{2-9,Y}={pagebreak=no},
    rowhead=1, leftsep=4pt, rightsep=4pt, rowsep=2.8pt
}
        CU-23 & Consultar una duda sobre las reglas \\
        Autor & Ibai Moya Aroz \\
        Requisitos asociados & \refRequisito{RF-4.2.1}, \refRequisito{RF-4.3}, \refRequisito{RF-4.4}, \refRequisito{RF-4.5}, \refRequisito{RF-4.6}, \refRequisito{RF-4.7}, \refRequisito{RF-4.8}, \refRequisito{RF-2.3.2} \\
        Actores & Usuario. \\
        Descripción & Permite aclarar una duda sobre las reglas de un juego y contrastar la respuesta con los manuales utilizados. \\
        Inicio & El usuario envía una pregunta sobre las reglas del juego. \\
        Precondiciones & El usuario ha iniciado sesión y el juego dispone de fuentes utilizables. Si continúa una conversación, esta pertenece a su cuenta. \\
        Acciones & \begin{pasosCU}
            \item El sistema utiliza la conversación seleccionada o crea una nueva asociada al juego.
            \item El sistema guarda la pregunta y muestra la respuesta como pendiente. Si es el primer mensaje, asigna un título a la conversación a partir de su contenido.
            \item El sistema genera la respuesta en segundo plano a partir de los manuales utilizables. Si hay mensajes anteriores, los tiene en cuenta para interpretar la pregunta.
            \item El sistema muestra la respuesta en español o inglés según el idioma detectado en el turno, junto con las referencias utilizadas.
            \item El usuario lee la respuesta y puede contrastarla consultando una página citada mediante \refCasoUso{CU-29}.
        \end{pasosCU} \\
        Postcondiciones & La pregunta y la respuesta generada quedan guardadas en la conversación. Si la generación falla, se conserva ese estado. Si esta es la primera actividad con el juego y el usuario no ha elegido antes seguirlo o dejar de seguirlo, el juego empieza a seguirse automáticamente. \\
        Alternativas & \begin{alternativasCU}
            \item Si no hay información suficiente para responder, el sistema lo indica.
            \item Si no quedan fuentes utilizables, no se admiten nuevos mensajes y el historial existente sigue disponible.
            \item Si falla el envío, el sistema muestra el error y recupera la pregunta para permitir otro intento.
            \item Si falla la generación, la pregunta queda guardada y la respuesta aparece como fallida.
        \end{alternativasCU} \\
\end{longtblr}

\clearpage

% CU-24
\begin{longtblr}[
    caption={\textbf{CU-24}. Consultar las conversaciones y su historial.},
    label={tab:cu:CU-24}
]{
    width=\linewidth,
    colspec={Q[l,wd=0.24\linewidth] X[l]},
    rows={valign=t},
    row{1}={bg=gray!18,font=\bfseries},
    column{1}={font=\bfseries},
    hline{1,Z}={0.8pt}, hline{2}={0.5pt},
    hborder{2-9,Y}={pagebreak=no},
    rowhead=1, leftsep=4pt, rightsep=4pt, rowsep=2.8pt
}
        CU-24 & Consultar las conversaciones y su historial \\
        Autor & Ibai Moya Aroz \\
        Requisitos asociados & \refRequisito{RF-4.2.1}, \refRequisito{RF-4.2.2}, \refRequisito{RF-4.6}, \refRequisito{RF-4.8} \\
        Actores & Usuario. \\
        Descripción & Permite recuperar las conversaciones propias de un juego para revisar sus mensajes y continuar las consultas anteriores. \\
        Inicio & El usuario accede a las conversaciones del juego. \\
        Precondiciones & El usuario ha iniciado sesión y el juego existe en la aplicación. \\
        Acciones & \begin{pasosCU}
            \item El sistema muestra la lista de conversaciones propias del juego.
            \item El usuario puede filtrar las conversaciones por su título para localizar la que desea consultar.
            \item Si abre una conversación, el sistema muestra sus mensajes, el estado de las respuestas y las referencias utilizadas. El usuario puede leer los mensajes y consultar sus referencias.
        \end{pasosCU} \\
        Postcondiciones & El usuario consulta sus conversaciones y los mensajes que desea revisar. Las operaciones opcionales pueden actualizar la lista o añadir una pregunta y su respuesta. Los mensajes anteriores se conservan, salvo que se elimine la conversación. \\
        Alternativas & \begin{alternativasCU}
            \item Si no hay conversaciones o el filtro no encuentra coincidencias, el sistema lo indica.
            \item Si se han retirado las fuentes utilizables, el historial se conserva y puede leerse, pero no se admiten nuevos mensajes.
            \item Si falla la carga de la lista o del historial, el sistema informa del error.
        \end{alternativasCU} \\
\end{longtblr}

\clearpage

% CU-25
\begin{longtblr}[
    caption={\textbf{CU-25}. Renombrar una conversación.},
    label={tab:cu:CU-25}
]{
    width=\linewidth,
    colspec={Q[l,wd=0.24\linewidth] X[l]},
    rows={valign=t},
    row{1}={bg=gray!18,font=\bfseries},
    column{1}={font=\bfseries},
    hline{1,Z}={0.8pt}, hline{2}={0.5pt},
    hborder{2-9,Y}={pagebreak=no},
    rowhead=1, leftsep=4pt, rightsep=4pt, rowsep=2.8pt
}
        CU-25 & Renombrar una conversación \\
        Autor & Ibai Moya Aroz \\
        Requisitos asociados & \refRequisito{RF-4.2.3} \\
        Actores & Usuario. \\
        Descripción & Permite identificar una conversación propia mediante un título elegido por el usuario. \\
        Inicio & El usuario elige la opción de renombrar una conversación. \\
        Precondiciones & El usuario ha iniciado sesión y la conversación pertenece a su cuenta. \\
        Acciones & \begin{pasosCU}
            \item El sistema muestra el título actual para que pueda modificarse.
            \item El usuario introduce el nuevo título y confirma el cambio.
            \item El sistema valida el título, lo guarda y actualiza la lista de conversaciones.
        \end{pasosCU} \\
        Postcondiciones & Tras guardar el cambio, la conversación aparece con el nuevo título y conserva sus mensajes. \\
        Alternativas & \begin{alternativasCU}
            \item Si el usuario cancela, se conserva el título anterior.
            \item Si el título está vacío o supera la longitud admitida, el sistema no permite guardarlo.
            \item Si no puede guardarse el cambio, el sistema informa del error.
        \end{alternativasCU} \\
\end{longtblr}

\clearpage

% CU-26
\begin{longtblr}[
    caption={\textbf{CU-26}. Eliminar una conversación.},
    label={tab:cu:CU-26}
]{
    width=\linewidth,
    colspec={Q[l,wd=0.24\linewidth] X[l]},
    rows={valign=t},
    row{1}={bg=gray!18,font=\bfseries},
    column{1}={font=\bfseries},
    hline{1,Z}={0.8pt}, hline{2}={0.5pt},
    hborder{2-9,Y}={pagebreak=no},
    rowhead=1, leftsep=4pt, rightsep=4pt, rowsep=2.8pt
}
        CU-26 & Eliminar una conversación \\
        Autor & Ibai Moya Aroz \\
        Requisitos asociados & \refRequisito{RF-4.2.4} \\
        Actores & Usuario. \\
        Descripción & Permite retirar una conversación propia y sus mensajes cuando el usuario ya no desea conservarlos. \\
        Inicio & El usuario solicita eliminar una de sus conversaciones. \\
        Precondiciones & El usuario ha iniciado sesión y la conversación pertenece a su cuenta. \\
        Acciones & \begin{pasosCU}
            \item El sistema pide confirmación antes de eliminar la conversación.
            \item El usuario confirma la eliminación.
            \item El sistema elimina la conversación y actualiza la lista.
        \end{pasosCU} \\
        Postcondiciones & Tras completar la eliminación, la conversación y sus mensajes dejan de ser accesibles. Las demás conversaciones se conservan. \\
        Alternativas & \begin{alternativasCU}
            \item Si el usuario cancela, se conservan la conversación y sus mensajes.
            \item Si la eliminación falla, el sistema informa del error y permite volver a intentarlo.
        \end{alternativasCU} \\
\end{longtblr}

\clearpage

% CU-29
\begin{longtblr}[
    caption={\textbf{CU-29}. Consultar una página citada.},
    label={tab:cu:CU-29}
]{
    width=\linewidth,
    colspec={Q[l,wd=0.24\linewidth] X[l]},
    rows={valign=t},
    row{1}={bg=gray!18,font=\bfseries},
    column{1}={font=\bfseries},
    hline{1,Z}={0.8pt}, hline{2}={0.5pt},
    hborder{2-9,Y}={pagebreak=no},
    rowhead=1, leftsep=4pt, rightsep=4pt, rowsep=2.8pt
}
        CU-29 & Consultar una página citada \\
        Autor & Ibai Moya Aroz \\
        Requisitos asociados & \refRequisito{RF-4.6}, \refRequisito{RF-3.6.1}, \refRequisito{RF-3.10} \\
        Actores & Usuario. \\
        Descripción & Permite comprobar una respuesta mediante la consulta de la página que se cita como fuente. \\
        Inicio & El usuario selecciona una referencia de una respuesta. \\
        Precondiciones & El usuario ha iniciado sesión y está leyendo una respuesta mediante \refCasoUso{CU-23} o \refCasoUso{CU-24}. La respuesta contiene una referencia a un manual propio disponible o a uno compartido que está activo. \\
        Acciones & \begin{pasosCU}
            \item El sistema abre el manual en la página indicada. Si es propio, permite su consulta mediante \refCasoUso{CU-18}; si es compartido y pertenece a otra cuenta, abre el visor en modo de solo lectura.
            \item El usuario consulta el texto y la imagen disponibles de la página para contrastar la respuesta.
        \end{pasosCU} \\
        Postcondiciones & El visor muestra la página citada, con el texto y la imagen disponibles. Si esa página ya no existe, muestra la primera página del manual. \\
        Alternativas & \begin{alternativasCU}
            \item Las referencias a manuales retirados o sin permiso de lectura se muestran como información, pero no permiten abrir su contenido.
            \item Si la página indicada ya no existe en el manual, el visor abre la primera página disponible.
            \item Si el manual deja de estar disponible antes de abrirlo, el sistema informa de que no puede consultarse.
        \end{alternativasCU} \\
\end{longtblr}

\clearpage

\subsubsection{Preferencias}

% CU-27
\begin{longtblr}[
    caption={\textbf{CU-27}. Seleccionar el idioma de la interfaz.},
    label={tab:cu:CU-27}
]{
    width=\linewidth,
    colspec={Q[l,wd=0.24\linewidth] X[l]},
    rows={valign=t},
    row{1}={bg=gray!18,font=\bfseries},
    column{1}={font=\bfseries},
    hline{1,Z}={0.8pt}, hline{2}={0.5pt},
    hborder{2-9,Y}={pagebreak=no},
    rowhead=1, leftsep=4pt, rightsep=4pt, rowsep=2.8pt
}
        CU-27 & Seleccionar el idioma de la interfaz \\
        Autor & Ibai Moya Aroz \\
        Requisitos asociados & \refRequisito{RF-5.1} \\
        Actores & Usuario. \\
        Descripción & El usuario adapta los textos de la interfaz al idioma en el que desea utilizar la aplicación. \\
        Inicio & El usuario decide cambiar el idioma de la interfaz. \\
        Precondiciones & El usuario puede acceder a la aplicación, con o sin una sesión iniciada. \\
        Acciones & \begin{pasosCU}
            \item El usuario selecciona español o inglés mediante el control de idioma.
            \item El sistema cambia los textos de la interfaz y guarda la selección en el dispositivo.
        \end{pasosCU} \\
        Postcondiciones & La interfaz utiliza el idioma elegido y conserva la selección al recargar la aplicación. \\
        Alternativas & \begin{alternativasCU}
            \item Durante la presentación inicial, el usuario permanece en el mismo paso después de cambiar el idioma.
            \item El idioma de las respuestas se determina a partir de la pregunta, según \refRequisito{RF-4.5}.
        \end{alternativasCU} \\
\end{longtblr}

\clearpage

% CU-28
\begin{longtblr}[
    caption={\textbf{CU-28}. Configurar la apariencia.},
    label={tab:cu:CU-28}
]{
    width=\linewidth,
    colspec={Q[l,wd=0.24\linewidth] X[l]},
    rows={valign=t},
    row{1}={bg=gray!18,font=\bfseries},
    column{1}={font=\bfseries},
    hline{1,Z}={0.8pt}, hline{2}={0.5pt},
    hborder{2-9,Y}={pagebreak=no},
    rowhead=1, leftsep=4pt, rightsep=4pt, rowsep=2.8pt
}
        CU-28 & Configurar la apariencia \\
        Autor & Ibai Moya Aroz \\
        Requisitos asociados & \refRequisito{RF-5.2} \\
        Actores & Usuario. \\
        Descripción & El usuario adapta el tema y el acento de color de la interfaz a sus preferencias. \\
        Inicio & El usuario abre los ajustes para cambiar la apariencia. \\
        Precondiciones & El usuario ha iniciado sesión y puede acceder a los ajustes de apariencia. \\
        Acciones & \begin{pasosCU}
            \item El usuario elige un tema claro, oscuro o automático, y puede cambiar el acento de color.
            \item El sistema aplica la configuración y la guarda en el dispositivo.
        \end{pasosCU} \\
        Postcondiciones & La interfaz muestra la apariencia elegida y conserva la configuración al recargar la aplicación. \\
        Alternativas & Si el usuario elige el tema automático, la interfaz sigue la preferencia del dispositivo. \\
\end{longtblr}

\clearpage

\subsubsection{Casos complementarios}

% CU-A1
\begin{longtblr}[
    caption={\textbf{CU-A1}. Consultar la ayuda y la presentación inicial.},
    label={tab:cu:CU-A1}
]{
    width=\linewidth,
    colspec={Q[l,wd=0.24\linewidth] X[l]},
    rows={valign=t},
    row{1}={bg=gray!18,font=\bfseries},
    column{1}={font=\bfseries},
    hline{1,Z}={0.8pt}, hline{2}={0.5pt},
    hborder{2-9,Y}={pagebreak=no},
    rowhead=1, leftsep=4pt, rightsep=4pt, rowsep=2.8pt
}
        CU-A1 & Consultar la ayuda y la presentación inicial \\
        Autor & Ibai Moya Aroz \\
        Requisitos asociados & \refRequisito{RNF-1}. \\
        Actores & Usuario. \\
        Descripción & El usuario consulta explicaciones sobre el funcionamiento y el uso de \textit{Manualito}. \\
        Inicio & El usuario accede a la presentación inicial o abre la sección de ayuda. \\
        Precondiciones & La presentación inicial es pública y se muestra si no se ha completado en el dispositivo. Para consultar la ayuda, el usuario debe haber iniciado sesión. \\
        Acciones & \begin{pasosCU}
            \item El sistema muestra las explicaciones disponibles en la presentación o en la ayuda.
            \item Si está en la presentación, el usuario recorre los pasos hasta la elección final de acceso o registro.
            \item Si está en la ayuda, consulta las explicaciones de uso y puede abrir las respuestas a las preguntas frecuentes.
        \end{pasosCU} \\
        Postcondiciones & El usuario ha consultado información sobre el uso de la aplicación. \\
        Alternativas & El usuario puede abandonar la presentación antes de completar sus pasos. \\
\end{longtblr}

\clearpage

% CU-A2
\begin{longtblr}[
    caption={\textbf{CU-A2}. Consultar la política de privacidad.},
    label={tab:cu:CU-A2}
]{
    width=\linewidth,
    colspec={Q[l,wd=0.24\linewidth] X[l]},
    rows={valign=t},
    row{1}={bg=gray!18,font=\bfseries},
    column{1}={font=\bfseries},
    hline{1,Z}={0.8pt}, hline{2}={0.5pt},
    hborder{2-9,Y}={pagebreak=no},
    rowhead=1, leftsep=4pt, rightsep=4pt, rowsep=2.8pt
}
        CU-A2 & Consultar la política de privacidad \\
        Autor & Ibai Moya Aroz \\
        Requisitos asociados & \refRequisito{RNF-6}. \\
        Actores & Usuario. \\
        Descripción & El usuario consulta la información publicada sobre el tratamiento de sus datos. \\
        Inicio & El usuario abre la política desde el registro, la presentación, los ajustes o su enlace directo. \\
        Precondiciones & No se requiere una sesión iniciada. \\
        Acciones & \begin{pasosCU}
            \item El sistema muestra la información publicada sobre el tratamiento de los datos.
            \item El usuario consulta el contenido de la política.
        \end{pasosCU} \\
        Postcondiciones & El usuario ha consultado la política de privacidad. \\
        Alternativas & \begin{alternativasCU}
            \item Si la abre desde el registro o la presentación, puede cerrar la consulta y continuar en el punto en el que se encontraba.
            \item La consulta de la política no sustituye la aceptación que se solicita durante el registro.
        \end{alternativasCU} \\
\end{longtblr}

\endgroup
